\section{Discussion and Conclusion}
In this work, we presented a theoretically grounded and optimization-free approach for efficient Language Gaussian Splatting. The proposed method shows state-of-the-art results with a single iteration over all views and avoids rendering and loss computation in the feature space.

Beyond the immediate performance gains, our approach opens up new possibilities for how we approach 3D scene understanding and manipulation. The ability to reason directly in language feature space, without compression, enables practical applications like object insertion and scene registration that were previously computationally prohibitive. Overall, we demonstrated that:
\begin{itemize}
    \item High-dimensional language features can directly be used in Gaussian splatting by utilizing a training-free optimization approach.
    \item Accurate modeling of the rendering process enables the backproject of highly relevant features from 2D to 3D.
    \item Directly projecting high-dimensional features allows us to use scene editing approaches previously limited to 3D Gaussian Splatting.
\end{itemize}
Based on the experiments and results presented in this work, we can answer the following questions about our work and its relation to existing works using feature compression:

\noindent\textbf{Is feature compression or spatial mapping necessary for Language Gaussian Splatting?} No. Avoiding the explicit rendering of language features and loss computation in the 2D space together with a training-free approach makes the direct use of high-dimensional language features feasible. 

\noindent\textbf{Should we stop using low-dimensional Language Gaussian Splatting?} No. Low-dimensional features are justified in certain scenarios despite constraints. Though our method enables the direct use of uncompressed features, low-dimensional alternatives can reduce memory usage, and accelerate 2D feature map rendering, and filter noise. However, it needs to be considered that scene-specific compression approaches may contradict generalization across scenes.

\noindent\textbf{Are there applications where low-dimensional language features are beneficial?} Yes. For scenarios where the domain of scenes is constrained, no scene editing is required, and storage and rendering cost exceeds the cost of training a feature compressor, low dimensional features are advantageous. Thus, the optimal feature resolution is strongly application-dependent and scales, especially in the domain of potential scenarios.

Overall, it is important to acknowledge that our approach is complementary to existing works in Language Gaussian Splatting and enables a new direction to train the 3D language representation. Due to this complementary nature, Occam's LGS can be combined with existing feature compression approaches, with the potential to speed up rendering and reduce storage costs. The right combination of approaches depends on the application and expected use case of the 3D representation, which is a direction we hope future work will explore.